\section{Further research questions}
\label{sec:questions}

\subsection{Higher-rank harmonic pencils and their moduli}
The rank-one case has the two exact forms
\eqref{eq:geometric-branch}--\eqref{eq:corner-branch}. Classify the
rational residue arrays for which the generic harmonic pencil has rank
at most two, then at most $\rho$. Hankel structure suggests a connection
with finite exponential-polynomial moment sequences and confluent
Vandermonde decompositions. The difficulty is the requirement that
\emph{all} coefficient matrices lie in one low-rank linear space, not
merely that each matrix separately has low rank. An effective necessary
and sufficient classification would extend the present finite residue
certificate beyond rank one.

\subsection{When does determinant degree equal harmonic rank?}
For a given rational weight, let $\delta$ be the total harmonic degree
of $\det(D_N+\mathscr C(\bm X))$. We proved $\delta\le\rho$.
Determine exact conditions for strict inequality, retaining the dependence
of $D_N$ on the same residue data as the pencil. A useful target is a
criterion in minors and discrete residues that avoids expanding the full
determinant. This would sharpen the formal polylogarithmic depth bound
from $\rho+1$ to $\delta+1$ before any functional reductions are used.

\subsection{Minimal depth in a specified functional relation system}
The generator \eqref{eq:explicitGF} contains $\Li_{2,2}(z,1)$, but
its coefficient could interact with functional equations in a broader
alphabet. Determine whether this generator reduces to classical
polylogarithms in a precisely declared differential or symbol calculus.
A nonzero obstruction would concern that formal system, not numerical
period independence. The problem is especially concrete because all
rational coefficients are explicit and the endpoint cancellation is
already known.

\subsection{Closed all-rank formulas for the finite exponential coefficients}
The determinants $U_r,V_r$ are finite at every rank, and $V_r(1,a)$
has the Gamma product \eqref{eq:residue-product}. Find short recurrences
or product formulas for $U_r(m,a)$ and $V_r(m,a)$ at fixed integer
$m\ge2$, or for their jets at $s=1$. The residue-product proof suggests
using finite discrete orthogonal polynomials at $m=1$ followed by
structured confluent perturbations. A successful result should improve
symbolic expression growth without introducing hidden infinite coordinates.

\subsection{Unsymmetrized mixed-order sums}
The polarized determinant evaluates the symmetrized sum with independent
orders and shifts. What additional coordinates are required for an
individual ordered assignment? For two particles, the sum and difference
under exchanging the two labels are natural first targets. The sum is
already evaluated by \eqref{eq:polarized-r2}; the antisymmetric component
may connect to multiple Hurwitz zeta or Lerch functions. One should seek
an identity on a genuine complex-parameter domain before differentiating
at integer values.

\subsection{Noninteger-order cutoff generators}
The spectral identity holds for arbitrary complex orders, but the rational
partial-fraction proof of polylogarithmic cutoff closure uses integer
orders. Determine a natural finite function class for the cutoff generator
at noninteger $s$. A likely candidate involves multiple Hurwitz--Lerch
functions with independently shifted denominators, but merely naming such
a function is not a reduction. The useful target is a finite transform
or differential system that specializes to Theorem~\ref{thm:explicitGF}
and preserves its initial conditions.

\subsection{More general annihilating stencils}
Replace the consecutive finite difference by
$\sum_j\alpha_j f(x+\lambda_j)$, with distinct rational shifts and
$\sum_j\alpha_j\lambda_j^k=0$ for $k<d$. Determine which stencils
still confine the moment special functions to one matrix direction.
When the shifts are integers, residue-polynomial divisibility is the
right language; on nonuniform grids, endpoint terms may produce several
Lerch coordinates. A classification should distinguish a genuinely
rank-one common image from a mere vanishing of the first asymptotic
coefficients.

\subsection{Higher Gauss resonances remain a distinct problem}
For the earlier Gauss ensemble, the resonances
$b-a-c=2r-2-M$ can have several divergent moment entries and overlapping
escape regimes. Our spectral deformation has only the single Hurwitz
pole and therefore does not resolve that general problem. Extend the
largest-index proof with a canonical subtraction of multiple escape
regions, and compare its cutoff polynomial with the Laurent coefficients
of the known Gamma product. The next two-particle case has a separate
incoming treatment; the remaining all-rank statement should retain that
provenance and its more complicated pole structure.

\subsection{Endpoint Euler primitives and weight drops}
Theorem~\ref{thm:depth} bounds the depth of every normalized Euler
primitive inside the disk. Determine which absolutely convergent endpoint
combinations of the integer-order finite-difference family reduce to
ordinary zeta values or products, and which require genuine multiple
zeta coordinates in a stated relation system. The cancellation in
\eqref{eq:AbelGF} is one exact example, not evidence for an unrestricted
closure claim. A systematic endpoint formula would connect the present
cutoff identities with the repository's harmonic-tail and cyclotomic
reduction programs.

\subsection{Formal verification with a small analytic interface}
The finite residue matrices, divisibility tests, polynomial annihilation,
and Cauchy--Binet normalization are suitable exact formalization targets.
They can be separated from three analytic interfaces: local normal
convergence of the spectral ensemble, the uniform largest-index remainder,
and the Hurwitz Laurent expansion. A proof-assistant implementation of
that separation would make subsequent repeated-pole identity generation
substantially more reliable. The supplied scripts are executable
certificates and diagnostics, not such a formalization.
